\loeskopf{Einheit 5 von 5 · Sachaufgaben (Nr. 38--42)}

\erg{38}{a) nur die positive \quad b) beide \quad c) nur die positive \quad
d) nur die positive \quad e) beide}

\erg{39}{a) $x_1 = 4$, $x_2 = -4$ \quad b) $x_1 = 14$, $x_2 = -14$ \quad
c) $x_1 = 20$, $x_2 = -20$ \quad d) $x_1 = 1{,}5$, $x_2 = -1{,}5$ \quad
e) $x^2 = 5x + 24$, also $x^2 - 5x - 24 = 0$; $x_1 = 8$, $x_2 = -3$ \quad
f) $x \cdot (x + 1) = 56$, also $x^2 + x - 56 = 0$; $x_1 = 7$, $x_2 = -8$ \quad
g) $x \cdot (x + 4) = 45$, also $x^2 + 4x - 45 = 0$; $x_1 = 5$, $x_2 = -9$}

\erg{40}{a) $x = $ Breite; $x \cdot 2x = 72$, also $2x^2 = 72$ und $x^2 = 36$;
$x_1 = 6$, $x_2 = -6$. Eine Länge ist nicht negativ. Breite $6\,$m, Länge $12\,$m.
\quad b) $x \cdot (x + 5) = 84$, also $x^2 + 5x - 84 = 0$; $x_1 = 7$, $x_2 = -12$.
$-12$ entfällt. Breite $7\,$m, Länge $12\,$m. \quad
c) $x = $ eine Seite, die andere ist $17 - x$ (halber Umfang $17\,$m);
$x \cdot (17 - x) = 60$, also $x^2 - 17x + 60 = 0$; $x_1 = 12$, $x_2 = 5$. Beide
Werte sind positiv und gehören zusammen: die Seiten sind $12\,$m und $5\,$m
(Probe: Umfang $34\,$m, Fläche $60\,\text{m}^2$).}

\erg{41}{Fehler: $-12$ ist keine Länge. Die Rechnung stimmt, aber die negative Lösung
muss im Sachzusammenhang ausgeschlossen werden; die zweite Seite ergibt sich aus
$x + 7$. Richtige Antwort: Breite $5\,$m, Länge $12\,$m. \quad
a) $x \cdot (x + 2) = 48$, also $x^2 + 2x - 48 = 0$; $x_1 = 6$, $x_2 = -8$; $-8$
entfällt. Breite $6\,$m, Länge $8\,$m.}

\erg{42}{a) Eine Seitenlänge ist eine Strecke und kann nicht negativ sein, deshalb
fällt die negative Lösung weg. Bei einem Zahlenrätsel ist jede Zahl erlaubt, also
zählen beide Lösungen -- es sei denn, die Frage schließt eine aus. \quad
b) Es gibt keine Zahl, die die Bedingung erfüllt: Die gesuchte Größe existiert unter
diesen Angaben nicht. Die Aufgabe hat keine Lösung, das ist kein Rechenfehler.}
